\documentclass[12pt]{article}
\usepackage[utf8]{inputenc}
\usepackage[T1]{fontenc}
\usepackage[margin=1in]{geometry}
\usepackage{setspace}
\usepackage{xcolor}
\usepackage[colorlinks=true, urlcolor=blue!60!black, linkcolor=blue!60!black]{hyperref}
\setlength{\parindent}{2em}
\setlength{\parskip}{1.25em}
\renewcommand{\baselinestretch}{1.0}
\begin{document}
\begin{center}
{\LARGE \textit{Dobbs v. Jackson Women's Health Organization}, 597 U.S. 215 (2022)\par}
\vspace{0.8em}
{\large Case Brief --- Four-Page Edition\par}
\end{center}
\vspace{1em}

\section{Memory Jogger}

The Court overruled \textit{Roe v. Wade} and \textit{Planned Parenthood v. Casey}, holding that the Constitution confers no right to abortion and returning the authority to regulate or prohibit abortion to the people and their elected representatives (\href{https://www.courtlistener.com/opinion/6481357/dobbs-v-jackson-womens-health-organization/}{CourtListener}).

\section{Detailed Case Facts}

Mississippi's 2018 Gestational Age Act prohibited abortion after 15 weeks' probable gestational age, with exceptions only for a medical emergency or a severe fetal abnormality and none for rape or incest. Jackson Women's Health Organization, the State's only licensed abortion clinic, which performed abortions up to 16 weeks, sued State Health Officer Thomas E. Dobbs. Under \textit{Roe} (1973) and \textit{Casey} (1992), a State could not prohibit abortion before viability (about 23 to 24 weeks) or impose an ``undue burden'' on pre-viability access. The case arrived after the appointments of Justices Gorsuch, Kavanaugh, and Barrett, while thirteen States had ``trigger laws'' ready; a draft of the majority opinion leaked in May 2022.

\section{Procedural History}

Judge Carlton W. Reeves (S.D. Miss.) granted summary judgment to the clinic and permanently enjoined the Act, since 15 weeks is ``prior to viability.'' The Fifth Circuit unanimously affirmed. The Supreme Court granted certiorari on whether ``all pre-viability prohibitions on elective abortions are unconstitutional,'' heard argument on December 1, 2021, and reversed and remanded on June 24, 2022.

\section{Judicial Votes}

\textbf{6--3} to uphold the 15-week ban; \textbf{5--4} to overrule \textit{Roe} and \textit{Casey}.
\begin{itemize}
\item \textbf{Majority:} Justice Alito, joined by Justices Thomas, Gorsuch, Kavanaugh, and Barrett.
\item \textbf{Concurrences:} Justices Thomas and Kavanaugh, each joining the majority.
\item \textbf{Concurring in the judgment:} Chief Justice Roberts.
\item \textbf{Joint dissent:} Justices Breyer, Sotomayor, and Kagan.
\end{itemize}

\section{Holding}

The Constitution does not confer a right to abortion; \textit{Roe} and \textit{Casey} are overruled. An unenumerated right is protected only if it is ``deeply rooted in this Nation's history and tradition'' and ``implicit in the concept of ordered liberty'' under \textit{Washington v. Glucksberg}. Abortion fails that test, and its regulation is not a sex-based classification, so abortion laws receive a ``strong presumption of validity'' and rational-basis review. Mississippi's Act survives. The assurance that nothing in the opinion casts doubt on non-abortion precedents is dicta.

\section{Analysis of Opinions}

\textbf{Majority (Alito).} After rejecting an Equal Protection theory under \textit{Geduldig v. Aiello}, Alito anchored due process ``liberty'' in history. From Bracton, Coke, Hale, and Blackstone through the nineteenth century, abortion was criminal at least after quickening, and by 1868 three-quarters of the States banned it at all stages. Abortion is ``critically different'' from \textit{Griswold}, \textit{Loving}, \textit{Lawrence}, and \textit{Obergefell} because it destroys ``fetal life.'' Five \textit{stare decisis} factors supported overruling: the nature of the error, \textit{Roe}'s ``exceptionally weak'' reasoning, the unworkable undue-burden test, distortion of other doctrines, and the absence of concrete reliance interests.

\textbf{Thomas (concurring).} Substantive due process is an oxymoron; ``in future cases, we should reconsider all of this Court's substantive due process precedents, including \textit{Griswold}, \textit{Lawrence}, and \textit{Obergefell}.''

\textbf{Kavanaugh (concurring).} The Constitution is ``neither pro-life nor pro-choice.'' Overruling \textit{Roe} does not threaten other precedents, and States may not bar interstate travel for abortion or punish past abortions retroactively.

\textbf{Roberts (concurring in the judgment).} He would discard the viability line and uphold the 15-week ban without overruling \textit{Roe} ``all the way down to the studs'': ``If it is not necessary to decide more to dispose of a case, then it is necessary not to decide more.''

\textbf{Joint dissent.} The ``1868 test'' freezes liberty at a time when women could not vote, consigning them to ``second-class citizenship.'' The majority's ``Scout's honor'' assurance cannot protect \textit{Griswold} or \textit{Obergefell}, which fail the same test, and nothing had changed except that ``the composition of this Court has changed.''

\section{Examples: Future Applications}

\begin{enumerate}
\item \textbf{Same-Side: Total Ban.} A ban from fertilization, with only a life exception, survives a federal due process challenge under rational-basis review.
\item \textbf{Same-Side: Telehealth Restriction.} A ban on prescribing abortion medication by telemedicine is upheld because a legislature could rationally think in-person supervision protects maternal health.
\item \textbf{Opposite-Side: Travel Ban.} A felony for traveling to another State for an abortion would likely fall under the right to interstate travel, as Justice Kavanaugh signaled.
\item \textbf{Opposite-Side: Contraception Ban.} A ban on contraception for married couples falls under \textit{Griswold}, which the majority declined to disturb.
\item \textbf{Fence-Sitter: IVF Embryos.} A law treating frozen embryos as ``unborn human beings'' pits \textit{Dobbs}'s history test against procreative liberty under \textit{Skinner} and \textit{Griswold}; IVF seeks to \textit{create} life.
\end{enumerate}

\section{Critique}

\textbf{Progressive critique.} Freezing ``liberty'' to 1868 perpetuates the disenfranchisement of women; historians dispute Alito's account of the common law; the \textit{Glucksberg} method cannot coherently spare \textit{Griswold}, \textit{Lawrence}, or \textit{Obergefell}; and dismissing reliance ignores generations of women's lived reliance, while overruling after personnel changes erodes legitimacy.

\textbf{Originalist critique.} Many originalists praise \textit{Dobbs} for restoring democratic self-government. Others fault it for keeping substantive due process instead of the Privileges or Immunities Clause, or for conflating original meaning with expected applications. Minimalists side with Roberts.

\section{Key Quotations}

\begin{quote}
``We hold that \textit{Roe} and \textit{Casey} must be overruled. The Constitution makes no reference to abortion, and no such right is implicitly protected by any constitutional provision, including the one on which the defenders of \textit{Roe} and \textit{Casey} now chiefly rely---the Due Process Clause of the Fourteenth Amendment.''
\end{quote}
\noindent --- Justice Alito, for the Court, 597 U.S. 215.

\begin{quote}
``\textit{Roe} was egregiously wrong from the start. Its reasoning was exceptionally weak, and the decision has had damaging consequences. And far from bringing about a national settlement of the abortion issue, \textit{Roe} and \textit{Casey} have enflamed debate and deepened division. It is time to heed the Constitution and return the issue of abortion to the people's elected representatives.''
\end{quote}
\noindent --- Justice Alito, for the Court, 597 U.S. 215.

\begin{quote}
``With sorrow---for this Court, but more, for the many millions of American women who have today lost a fundamental constitutional protection---we dissent.''
\end{quote}
\noindent --- Justices Breyer, Sotomayor, and Kagan, dissenting, 597 U.S. 215.

\end{document}
